\documentclass[8pt]{beamer}

\usetheme{Darmstadt}
\usecolortheme{dove}
\usepackage[small,center]{caption}
\usepackage{times}
\usefonttheme{structurebold}
\usepackage[english]{babel}
\usepackage{pgf,pgfarrows,pgfnodes,pgfautomata,pgfheaps}
\usepackage{amsmath,amssymb}
\usepackage{amsxtra}

\usepackage{caption}
\usepackage{tikz}
\usetikzlibrary{shapes.misc}

\newcommand{\leftd}[1]{{\color{red} \bar{#1}}}
\newcommand{\interface}[2]{{\color{blue}{#1}_I(#2)}}
\newcommand{\leftdd}[2]{{\color{red} \bar{#1}(\bar{#2})}}
\newcommand{\leftFourier}[1]{{\color{red} \hat{#1}}}
\newcommand{\leftFourierTwo}[2]{{\color{red} \hat{#1}(\hat{#2})}}
\newcommand{\half}{\frac{1}{2}}
\newcommand{\divergence}{\mathrm{div}}

\newcommand{\I}{I}

\newcommand{\vst}{\vskip 0.2cm\noindent}
\newcommand{\vso}{\vskip 0.1cm\noindent}
\newcommand{\ba}{\begin{array}}
\newcommand{\ea}{\end{array}}
\newcommand{\be}{\begin{equation}}
\newcommand{\ee}{\end{equation}}
\newcommand{\matl}{M{\scriptsize ATLAB }}
\newcommand{\Div}[1]{\nabla\cdot #1}
\newcommand{\DIV}[1]{\nabla_{\mathbf{X}} \cdot #1}
\newcommand{\kk}{\mathcal{K}}

\newcommand{\etal}{\emph{et al}.}

\newcommand{\BB}{\mathbb{B}}
\newcommand{\CC}{\mathbb{C}}
\newcommand{\FF}{\mathbb{F}}
\newcommand{\EE}{\mathbb{E}}
\newcommand{\PP}{\mathbb{P}}
\newcommand{\MM}{\mathbb{M}}
\newcommand{\FFb}{\overline{\FF}}

\newcommand{\ub}{\boldsymbol{u}}
\newcommand{\ubf}{\boldsymbol{u}^{\text{f}}}
\newcommand{\ubs}{\boldsymbol{u}^{\text{s}}}
\newcommand{\vb}{\boldsymbol{v}}
\newcommand{\wb}{\boldsymbol{w}}
\newcommand{\yb}{\boldsymbol{y}}
\newcommand{\Chib}{\boldsymbol{\chi}}
\newcommand{\cb}{\boldsymbol{\chi}}
\newcommand{\xb}{\boldsymbol{x}}
\newcommand{\Xb}{\boldsymbol{X}}
\newcommand{\fb}{\boldsymbol{f}}
\newcommand{\Fb}{\boldsymbol{F}}
\newcommand{\Ub}{\boldsymbol{U}}
\newcommand{\UbVecIB}{\vec{\Ub}\vphantom{\Ub}^{\text{IB}}}
\newcommand{\Vb}{\boldsymbol{V}}
\newcommand{\Lb}{\boldsymbol{L}}
\newcommand{\LbVecIB}{\vec{\Lb}\vphantom{\Lb}^{\text{IB}}}
\newcommand{\Nb}{\boldsymbol{N}}
\newcommand{\nb}{\boldsymbol{n}}
\newcommand{\hb}{\boldsymbol{h}}
\newcommand{\ab}{\boldsymbol{a}}
\newcommand{\Ab}{\boldsymbol{A}}
\newcommand{\pb}{\boldsymbol{\phi}}

\newcommand{\cauchy}{\bbsigma}
\newcommand{\cauchyf}{\bbsigma^{\text{f}}}
\newcommand{\cauchyv}{\bbsigma^{\text{v}}}
\newcommand{\cauchys}{\bbsigma^{\text{e}}}
\newcommand{\cauchysolid}{\bbsigma^{\text{s}}}
\newcommand{\ztensor}{\mathbb{0}}
\newcommand{\pstab}{\pi_{\text{stab}}}
\newcommand{\pf}{\pi^{\text{f}}}
\newcommand{\ps}{\pi^{\text{s}}}
\newcommand{\pphys}{p}
\newcommand{\peul}{\pi}
\newcommand{\mus}{\mu^{\text{s}}}
\newcommand{\muf}{\mu^{\text{f}}}
\newcommand{\rhos}{\rho^{\text{s}}}
\newcommand{\rhof}{\rho^{\text{f}}}
\newcommand{\cbs}{\Chib^{\text{s}}}
\newcommand{\ds}{\boldsymbol{d}^{\text{s}}}

\newcommand{\PPs}{\PP^{\text{e}}}
\newcommand{\PPsm}{\bar{\PP}^{\text{e}}}

\newcommand{\soliddom}{{\Omega^{\text{s}}_t}}
\newcommand{\soliddomO}{{\Omega^{\text{s}}_0}}
\newcommand{\fluiddom}{\Omega^{\text{f}}_t}
\newcommand{\fluiddomO}{\Omega^{\text{f}}_0}
\newcommand{\fsinterface}{\Gamma^{\text{fs}}_t}
\newcommand{\fsinterfaceO}{\Gamma^{\text{fs}}_0}
\newcommand{\subsolid}{\omega^{\text{s}}_t}
\newcommand{\subsolidO}{\omega^{\text{s}}_0}

\newcommand{\Wiso}{W^{\text{iso}}}
\newcommand{\Wb}{\overline{W}}

\newcommand{\PPb}{\bar{\PP}^\text{e}}

\newcommand{\cNHB}{\bar{\cauchy}_{\text{NH}}}
\newcommand{\dev}{\text{dev}}
\newcommand{\DEV}{\text{DEV}}
\newcommand{\efac}{E_\text{FAC}}
\newcommand{\mfac}{M_\text{FAC}}
\newcommand{\tr}{\text{tr}}
\newcommand{\nus}{\nu_\text{stab}}
\newcommand{\kappas}{\kappa_\text{stab}}
%\newcommand{\mus}{\mu_\text{s}}
%\newcommand{\mut}{\mu_\text{T}}
\newcommand{\mul}{\mu_\text{L}}
\newcommand{\Gt}{G_\text{T}}
\newcommand{\Gl}{G_\text{L}}
\newcommand{\El}{E_\text{L}}

\newcommand{\vol}{\text{vol}}
\newcommand{\nvol}{\mathcal{N}}
\newcommand{\supp}{\text{supp}}
\newcommand{\mm}{\mathcal{M}}
\newcommand{\mmd}{\mathcal{M}^{\text{1D}}}
\newcommand{\xib}{\boldsymbol{\xi}}
\newcommand{\Pone}{\mathcal{P}^1}
\newcommand{\Ptwo}{\mathcal{P}^2}
\newcommand{\Qone}{\mathcal{Q}^1}
\newcommand{\Qtwo}{\mathcal{Q}^2}
\newcommand{\Bb}{\boldsymbol{B}}
\newcommand{\order}{\mathcal{O}}

\newcommand{\tria}{{\ensuremath{\mathcal{T}^h}}}
\newcommand{\euleriandx}{{\ensuremath{\Delta x}}}
\newcommand{\lagrangiandx}{{\ensuremath{\Delta X}}}

\newcommand{\dAb}{\, \mathrm{d}\Ab}
\newcommand{\dXb}{\, \mathrm{d}\Xb}

\newcommand{\dxb}{\, \mathrm{d}\xb}

% fix placement of superscripts on characters with tildes
\newcommand{\tilden}[2]{\ensuremath{\tilde{#1}\vphantom{#1}^{#2}}}

\newcommand{\xface}{{i, j + \half, k + \half}}
\newcommand{\yface}{{i + \half, j, k + \half}}
\newcommand{\zface}{{i + \half, j + \half, k}}

% same with transposes and vecs
\newcommand{\vecT}[1]{\ensuremath{\vec{#1}\vphantom{#1}^T}}
% getting the arrow to appear nicely over the 1 is surprisingly tricky
\newcommand{\vecOneT}{%
  \ensuremath{%
    \vec{\vphantom{\boldsymbol{1}}\hphantom{-}}% Put an arrow over a box of about the right shape
    \llap{\ensuremath{\boldsymbol{1}}}% place the symbol to the left of the current location of the cursor
    ^T%
  }}

\newcommand*{\vcenterarrow}{\vcenter{\hbox{$\Longrightarrow$}}}

\DeclareMathOperator{\hyphen}{-}

\definecolor{Carolinablue}{RGB}{75, 156, 211}
\definecolor{ballblue}{rgb}{0.13, 0.67, 0.8}
\definecolor{lightgray}{rgb}{0.83, 0.83, 0.83}
\setbeamercolor{block title}{bg=lightgray,fg=Carolinablue}
\setbeamercolor{block body}{bg=white,fg=black}
\setbeamercovered{dynamic}
\setbeamercolor*{item}{fg=Carolinablue}

\captionsetup[subfigure]{labelformat=empty}
\captionsetup[figure]{labelformat=empty}
\setbeamertemplate{navigation symbols}{}
\setbeamertemplate{footline}[frame number]
\begin{document}

% tikz stuff
\tikzset{cross/.style={cross out, draw=black, minimum size=2*(#1-\pgflinewidth),
inner sep=0pt, outer sep=0pt},
%default radius will be 1pt.
cross/.default={2.5pt}}



\frame{
  \title{\Large deal.II, IBAMR, and Cardinal: Scalable Software for Difficult Problems}

\author{{\Large David Wells \\\vspace{0.1in} University of North Carolina}    \\
\vspace{0.2in} {Cardiovascular Modeling and Simulation Lab}}

\date{October 2, 2024\\{}}

\begin{figure}[h]
\centering
\includegraphics[width=1.5in]{UNC_logo_542.eps}
 \end{figure}%

\vspace{-0.2in}
\titlepage
}

\section{Overview}
\begin{frame}
    \frametitle{Overview}
    {\large About Me}
    \begin{itemize}
    \item[$\blacksquare$] PhD in applied mathematics, 2015
    \item[$\blacksquare$] postdoc at RPI, 2015 - 2018
    \item[$\blacksquare$] postdoc at UNC, 2018 - 2023
    \item[$\blacksquare$] research scientist at UNC, 2023 - present
    \end{itemize}

    \vfill

    {\large Why?}
    \begin{itemize}
    \item[$\blacksquare$] Help students on their journeys in scientific software
    \item[$\blacksquare$] Contribute to the greater applied mathematics community
      by maintaining software, contributing patches to upstream projects
    \item[$\blacksquare$] Solve new and interesting problems with cutting-edge tools
    \end{itemize}

    \vfill

    {\large Topics}
    \begin{itemize}
    \item[$\blacksquare$] deal.II
    \item[$\blacksquare$] IBAMR
    \item[$\blacksquare$] Cardinal
    \end{itemize}
\end{frame}

\section{deal.II}
\begin{frame}
    \frametitle{deal.II: Overview}
    \begin{figure}
        \centering
        \includegraphics[width=1.5in]{step-32.3d-sphere.partition.cropped.png}
    \end{figure}
    \begin{itemize}
        \item[$\blacksquare$] Large C++ library implementing the adaptive finite element method
        \item[$\blacksquare$] Extraordinarily well-documented with an active 
                             community, classes, videos, examples, etc.
        \item[$\blacksquare$] Scalable to largest machines today (in use at ORNL
          and LANL, hetereogeneous compute is WIP)
        \item[$\blacksquare$] Successfully used by students on laptops to solve dissertation-level problems
        \item[$\blacksquare$] Successful industry applications (Siemens, Italian
          shipbuilding, Canadian lumber industry)
        \item[$\blacksquare$] Does not reinvent the wheel: uses Kokkos, MPI, PETSc, Trilinos, p4est
    \end{itemize}
\end{frame}

\begin{frame}
    \frametitle{deal.II: Community Impact}
    \begin{figure}
        \centering
        \includegraphics[width=4in]{publications.png}
    \end{figure}
    \begin{itemize}
        \item[$\blacksquare$] I've been a principal developer since 2015 (one of fourteen)
        \item[$\blacksquare$] Large, quantifiable impact on science
        \item[$\blacksquare$] S-tier finite element library
        \item[$\blacksquare$] If you need to solve a PDE, deal.II can help
    \end{itemize}
\end{frame}

\begin{frame}
    \frametitle{deal.II Showcase: plate tectonics}
    \begin{figure}
        \centering
        \includegraphics[width=2in]{aspect.png}
    \end{figure}
    \begin{itemize}
        \item[$\blacksquare$] ASPECT: solve Stokes system at extreme scale (ten trillion DoFs) on Frontera (Gassmoeller et al)
        \item[$\blacksquare$] Picture: \emph{Changes in core-mantle boundary heat flux
          patterns throughout the supercontinent cycle}, Dannberg et al: how does heat flux inside the earth impact our magnetic field?
        \item[$\blacksquare$] Matrix-free geometric multigrid, adaptive FEM,
          scalable linear solvers and preconditioners
    \end{itemize}
\end{frame}

\begin{frame}
    \frametitle{deal.II Showcase: hypersonic flows}
    \begin{figure}
        \centering
        \includegraphics[width=2in]{ryujin.png}
    \end{figure}
    \begin{itemize}
        \item[$\blacksquare$] ryujin: convex limiting, graph viscosity methods
          (Guermond and Popov) for hyperbolic conservation laws (Maier et al)
        \item[$\blacksquare$] Picture: Compressible Euler equation with
          Jones-Wilkins-Lee equation of state (Maier et al)
        \item[$\blacksquare$] Not just for EOS-based hypersonic flows: also used
          to model flooding in Texas with shallow-water equations
    \end{itemize}
\end{frame}

\section{IBAMR}
\begin{frame}
    \frametitle{IBAMR: Overview}
    \begin{figure}
        \centering
        \includegraphics[width=2in]{cardiac.png}
    \end{figure}
    \begin{itemize}
        \item[$\blacksquare$] Large C++ library with both finite differences (FDM) and finite elements (FEM)
        \item[$\blacksquare$] Scalable 2d and 3d implementations of common IB methods
        \item[$\blacksquare$] Major project, lead by Boyce Griffith here at UNC, to solve hard FSI problems,
             complex fluids, multiphase flows (Bhalla et al at SDSU)
        \item[$\blacksquare$] I have been a co-maintainer since 2018
    \end{itemize}
\end{frame}

\begin{frame}
    \frametitle{IBAMR Showcase: Numerical Wave Tank}
    \begin{figure}
        \centering
        \includegraphics[width=2in]{wavetank.png}
    \end{figure}
    \begin{itemize}
        \item[$\blacksquare$] By the group at SDSU lead by Bhalla
        \item[$\blacksquare$] Picture: \emph{A model predictive control (MPC)-integrated
        multiphase immersed boundary (IB) framework for simulating wave energy
        converters (WECs)}, Khedkar et al
        \item[$\blacksquare$] Model solids, liquids, and gasses efficiently to solve a real problem
    \end{itemize}
\end{frame}

\begin{frame}
    \frametitle{IBAMR Showcase: Transarterial Valve Replacement}
    \begin{figure}
        \centering
        \includegraphics[width=3in]{tavr.png}
    \end{figure}
    \begin{itemize}
        \item[$\blacksquare$] What happens when we replace someone's tricuspid valve?
        \item[$\blacksquare$] Picture: \emph{Patient–Specific Immersed Finite Element–Difference Model
of Transcatheter Aortic Valve Replacement}, Brown et al: stent and valve and von Mises stresses
        \item[$\blacksquare$] Help surgeons understand how these devices work with detailed simulations
    \end{itemize}
\end{frame}

\section{Cardinal}
\begin{frame}
    \frametitle{Cardinal: IBAMR plus deal.II}
    \begin{figure}
        \centering
        \includegraphics[width=1.5in]{metermesh.png}
    \end{figure}
    \begin{itemize}
        \item[$\blacksquare$] Unified interface to both fluid-structure
            interaction (FSI) and electrophysiology (EP) models
        \item[$\blacksquare$] Picture: codim 1 meter mesh measuring flow through
            a tricuspid valve
        \item[$\blacksquare$] Detailed models (HGO, O'Hara-Rudy) practitioners care about
        \item[$\blacksquare$] Group here with Wells, Abdala, Davey, Griffith, and others
        \item[$\blacksquare$] Builds on deal.II's scalable FE infrastructure and 
            IBAMR's IB and incompressible Navier-Stokes (INS) solvers
    \end{itemize}
\end{frame}

\begin{frame}
    \frametitle{Cardinal: IBAMR plus deal.II}
    \begin{figure}
        \centering
        \includegraphics[width=2in]{mitral.png}
    \end{figure}
    \begin{itemize}
        \item[$\blacksquare$] Picture: flow around mitral valve and chordae
        \item[$\blacksquare$] Project 1: fully distributed data structures for IB (Wells)
        \item[$\blacksquare$] Project 2: better patch-based load balancers (Davey, Wells)
        \item[$\blacksquare$] Project 3: right heart model (Davey, collaboration with UC Irvine)
        \item[$\blacksquare$] Project 4: fully coupled EP-FSI (Abdala)
        \item[$\blacksquare$] Project 5: distributed Lagrange multipliers (TBA)
    \end{itemize}
\end{frame}

\end{document}
